\begin{textsample}{POS Dim 1 – ba – Score 74.00 – ba/ba\_ef\_15.txt}  \label{ex:f1_pos_102}
8.3.
ÁREA DE MATEMÁTICA – TEXTO \textbf{INTRODUTÓRIO} O Currículo Baiano da área de Matemática propõe a ampliação e o \textbf{aprofundamento} das aprendizagens essenciais desenvolvidas até o 9º ano do Ensino Fundamental.
A presente proposta considera que a sociedade \textbf{contemporânea}, ao realizar ações das mais \textbf{simples} até aquelas que envolvem conceitos científicos e tecnológicos, utiliza conhecimentos matemáticos que vão sendo construídos historicamente pelas necessidades diárias dos indivíduos.
Nessa perspectiva, para que a escola acompanhe a história da civilização, ou seja, o processo de desenvolvimento humano que se encontra \textbf{ancorado} no contexto da resolução de situações-problema, deve -se conceber uma nova dinâmica para a \textbf{mobilização} de saberes matemáticos intrinsecamente ligados a uma realidade sociocultural.
Para \textbf{tanto}, colocam -se em jogo, de modo mais inter-relacionado, os conhecimentos já explorados na etapa anterior, de modo a possibilitar que os estudantes construam uma visão mais integrada da Matemática, ainda na perspectiva de sua aplicação à realidade.
Assim, cada escola deve ser suficientemente flexível para contemplar os estudantes de diferentes níveis de habilidade e deve espelhar -se em suas necessidades – entre estas \textbf{figuram} experiências matemáticas significativas e interessantes sobre outras áreas de aprendizagem.
Além disso, deve oportunizar a \textbf{compreensão} da necessidade de continuarem estudando Matemática além dos muros da escola ; e uma formação como sujeitos alfabetizados matematicamente, capazes de fazer uso social das habilidades e \textbf{competências} construídas no Ensino Fundamental.
Pensar na contribuição da Matemática na construção do Currículo por \textbf{Competência} tem sido um desafio para os educadores do mundo \textbf{contemporâneo}.
Desenvolvida a partir das civilizações mediterrâneas, desde egípcios, babilônicos, hebreus, gregos e romanos, e que, a partir do \textbf{século} XVI, passou a todas as regiões do planeta, a Matemática não nasceu como ela se apresenta hoje.
Ela nasceu do esforço de lidar com questões do dia a dia, formalizando as ideias matemáticas, das práticas sociais, da relação do homem como seu meio e da necessidade de resolver \textbf{problemas} postos em seu contexto de vida, valorizando o conhecimento preexistente ao ingressar na escola.
O currículo ora proposto visa proporcionar ao estudante uma Educação Matemática \textbf{tanto} adequada, do ponto de \textbf{vista} escolar, quanto socialmente \textbf{relevante}.
A importância da Matemática deve ser vista como processo de construção de conhecimento, favorecido mediante a estimulação da investigação e participação dos \textbf{alunos}, o qual possa estar relacionado aos avanços tecnológicos, social e cultural da história da \textbf{humanidade}, seja por sua grande aplicação na sociedade \textbf{contemporânea}, seja pelas suas potencialidades na formação de cidadãos \textbf{críticos}, cientes de suas responsabilidades.
A Matemática e a língua materna – entendida aqui como a primeira língua que se aprende – têm sido as \textbf{disciplinas} básicas na constituição dos currículos escolares, em todas as épocas e culturas, havendo um razoável \textbf{consenso} ao fato de que, sem o desenvolvimento adequado de tal eixo linguístico/lógico-matemático, a formação pessoal não se completa.
Portanto, os documentos que vislumbram seu funcionamento devem contemplar o presente e olhar para o futuro.
E, nesse olhar para o futuro, \textbf{espera} -se que a escola, paralelamente ao ensino dos conteúdos escolares, esteja apta para habilitar os jovens estudantes com \textbf{competências} que lhes permitam trabalhar em equipe e intervir de forma \textbf{crítica}, consciente e autônoma.
Ou seja, \textbf{exige} -se a elaboração de um Currículo por \textbf{Competência}, aqui proposto, que possa oferecer uma educação de qualidade para todos : uma educação plural, democrática, inclusiva e hábil na construção de uma sociedade \textbf{baseada} em valores de cooperação, parceria e solidariedade.
No entanto, esse referencial só terá significado se todos os \textbf{atores} \textbf{implicados} no processo educativo apropriarem -se dele.
É \textbf{imprescindível} destacar que as \textbf{competências} gerais da Educação Básica, apresentadas a seguir, \textbf{inter-relacionam} -se e desdobram -se no tratamento didático proposto para as três etapas da Educação ( Educação Infantil, Ensino Fundamental e Ensino Médio ).
Elas estão organizadas e relacionam -se, de modo progressivo, para culminar no desenvolvimento das dez \textbf{competências} gerais apontadas pela BNCC, que são : 1.
Valorizar e utilizar os conhecimentos historicamente construídos sobre o mundo físico, social, cultural e \textbf{digital} para entender e explicar a realidade, continuar aprendendo e colaborar para a construção de uma sociedade justa, democrática e inclusiva. 2.
Exercitar a curiosidade intelectual e recorrer à \textbf{abordagem} própria das \textbf{ciências}, incluindo a investigação, a reflexão, a \textbf{análise} \textbf{crítica}, a imaginação e a criatividade, para investigar causas, elaborar e testar hipóteses, formular e resolver \textbf{problemas} e criar soluções ( inclusive tecnológicas ) com base nos conhecimentos das diferentes áreas. 3.
Valorizar e fruir as diversas manifestações artísticas e culturais, das locais às mundiais, e também participar de práticas diversificadas da produção \textbf{artístico-cultural}. 4.
Utilizar diferentes linguagens – verbal ( oral ou \textbf{visual-motora}, como \textbf{Libras}, e escrita ), corporal, visual, sonora e \textbf{digital} –, bem como conhecimentos das linguagens artística, matemática e científica, para se expressar e partilhar informações, experiências, ideias e sentimentos em diferentes contextos e produzir sentidos que levem ao entendimento mútuo. 5.
Compreender, utilizar e criar \textbf{tecnologias} \textbf{digitais} de informação e comunicação de forma \textbf{crítica}, significativa, \textbf{reflexiva} e ética nas diversas práticas sociais ( incluindo as escolares ) para se comunicar, acessar e disseminar informações, produzir conhecimentos, resolver \textbf{problemas} e exercer protagonismo e \textbf{autoria} na vida pessoal e coletiva. 6.
Valorizar a diversidade de saberes e vivências culturais e apropriar -se de conhecimentos e experiências que lhe possibilitem entender as relações próprias do mundo do trabalho e fazer escolhas alinhadas ao exercício da cidadania e ao seu projeto de vida, com liberdade, autonomia, consciência \textbf{crítica} e responsabilidade. 7. \textbf{Argumentar} com base em fatos, dados e informações confiáveis, para formular, negociar e defender ideias, pontos de \textbf{vista} e decisões comuns que respeitem e promovam os direitos humanos, a consciência \textbf{socioambiental} e o consumo responsável em âmbito local, regional e global, com \textbf{posicionamento} ético em relação ao cuidado de si mesmo, dos outros e do planeta. 8.
Conhecer -se, apreciar -se e cuidar de sua saúde física e emocional, compreendendo -se na diversidade humana e reconhecendo suas emoções e as dos outros, com autocrítica e \textbf{capacidade} para lidar com elas. 9.
Exercitar a empatia, o diálogo, a resolução de conflitos e a cooperação, fazendo -se respeitar e promovendo o respeito ao outro e aos direitos humanos, com acolhimento e valorização da diversidade de indivíduos e de grupos sociais, seus saberes, identidades, culturas e potencialidades, sem preconceitos de qualquer natureza. 10.
Agir pessoal e coletivamente com autonomia, responsabilidade, flexibilidade, resiliência e determinação, tomando decisões com base em princípios éticos, democráticos, inclusivos, sustentáveis e solidários.
É importante destacar que essas \textbf{competências} articulam -se na construção de conhecimentos, no desenvolvimento de habilidades e na formação de atitudes e valores, nos termos da LDB.
A escola é um espaço onde \textbf{inúmeras} pessoas interagem com intencionalidades e responsabilidades definidas.
Essa organização constitui um ambiente de aprendizagem, cuja atmosfera pode propiciar uma vivência do que queremos como sociedade : um espaço de igualdade, acolhedor da diversidade, onde o conhecimento e as relações interpessoais favorecem a inserção e um olhar amplo para o que acontece no mundo.
Para que haja avanço e ampliação dos esquemas mentais, da rede de conhecimentos, na qual essa organização do ambiente de aprendizagem possa fluir naturalmente, é papel do(da) professor(a) propor atividades que gerem conflitos para que os \textbf{alunos} compreendam as ligações entre os conteúdos, valorizando sempre a construção do conhecimento significativo.
Dentre as \textbf{inúmeras} atividades, a Matemática é privilegiada com brincadeiras, jogos e \textbf{problemas}, além de muitas outras atividades que \textbf{auxiliam} no desenvolvimento integral da criança, possibilitando observar, refletir, interpretar, levantar hipóteses, procurar e descobrir explicações ou soluções, expressar ideias e sentimentos – são desafios essenciais a serem propostos no processo educativo.
A transição da Educação Infantil para o Ensino Fundamental requer uma atenção singular e especial, visto que as aprendizagens e o desenvolvimento das crianças têm como eixos estruturantes as interações e a brincadeira, assegurando -lhes os direitos de conviver, brincar, participar, explorar, expressar -se e conhecer -se.
Nessa proposta, o DCRB traz a alfabetização e o \textbf{letramento} matemático em todos os momentos das unidades temáticas como eixos estruturantes e como aspectos \textbf{transversais} no Ensino Fundamental.
A alfabetização matemática é o processo de aprendizagem do sistema da escrita numérica na qual se desenvolve a habilidade de ler e escrever matematicamente.
O \textbf{letramento} matemático constitui -se no processo de apropriação do uso competente da leitura e da escrita nas práticas sociais.
Esse \textbf{letramento} expressa uma ideia social de leitura e escrita, na qual os sujeitos compreendem os números no seu cotidiano, por meio da interação com o meio e com os outros e das experiências das situações vivenciadas.
Esse é o \textbf{caminho} que o DCRB pretende potencializar junto com a BNCC, que traz na sua proposta o \textbf{foco} do que precisa ser desenvolvido no \textbf{aluno}, para que o conhecimento matemático seja uma ferramenta para ler, compreender e transformar a realidade.
Assim, propõe -se um compromisso muito forte com o \textbf{letramento} matemático, definido como as \textbf{competências} e habilidades de raciocinar, representar, comunicar e \textbf{argumentar} matematicamente, de modo a favorecer o estabelecimento de conjecturas, a formulação e a resolução de \textbf{problemas} em uma variedade de contextos, utilizando conceitos, procedimentos, fatos e ferramentas matemáticas.
O desenvolvimento dessas habilidades está intrinsecamente relacionado a algumas formas de organização da aprendizagem matemática, com base na \textbf{análise} de situações da vida cotidiana, de outras áreas do conhecimento e da própria Matemática.
Os processos matemáticos de resolução de \textbf{problemas}, de investigação, de desenvolvimento de projetos e da modelagem podem ser citados como formas privilegiadas da atividade matemática, motivo pelo qual são, ao mesmo tempo, objeto e estratégia para a aprendizagem ao longo de todo o Ensino Fundamental dos Anos Iniciais ou \textbf{Finais}.
Tais processos são potencialmente ricos para o desenvolvimento de \textbf{competências} fundamentais para a alfabetização e o \textbf{letramento} matemático.
Para o desenvolvimento de \textbf{competências} e habilidades matemáticas, como resolver \textbf{problemas} não somente escolares, mas também de práticas cotidianas e sociais, tais como : ler gráficos e tabelas, interpretar contas de água, luz, telefone, entre outras ações que \textbf{dependem} de conhecimentos relacionados aos diferentes usos socioculturais da Matemática, e propor um currículo que seja vivo no sentido de valorizar, principalmente, os processos de ensino e aprendizagem de Matemática, que acontecem em diversas situações e em múltiplos ambientes, desde o convívio em casa até os grupos sociais – escola, parque, igrejas etc.
Sobre isso, a BNCC ( BRASIL, 2017, p. 118 ) aborda que o currículo de Matemática deve \textbf{aproximar} as temáticas de Matemática ao universo da cultura, das \textbf{contextualizações} [... ].
Não podemos negar a relevância da escola a partir das atividades em sala de \textbf{aula} realizadas pelos professores e \textbf{alunos}, pois é na escola que acontecem as interações/mediações que possibilitam a consolidação do aprendizado, possibilitando o desenvolvimento do sujeito \textbf{aprendiz} como capaz de raciocinar, analisar, deduzir, criar, resolver situações e buscar estratégias inovadoras.
Assim, no \textbf{intuito} de acontecer a construção de conhecimento do \textbf{aluno} a partir do fazer matemático, é necessário gerar em sala de \textbf{aula} uma prática de produção em que os \textbf{alunos} se apropriem não só dos saberes, mas também vivenciem uma \textbf{abordagem} efetiva dos saberes matemáticos presentes no contexto cultural, em especial da cultura baiana, recuperando e valorizando as tradições culturais, festividades, história de cada região do nosso estado, além dos patrimônios histórico e artístico encontrados em todos os entornos das escolas baianas.
Nos dias \textbf{atuais}, há uma grande necessidade de que os(as) professores(as) desenvolvam \textbf{competências} profissionais para preparar os \textbf{alunos} em uma formação crítico-social.
Para \textbf{tanto}, é preciso substituir as formas tradicionais de ensino por aprendizagens ativas que tornem o \textbf{aluno} protagonista do seu próprio processo de \textbf{ensino-aprendizagem}, valendo -se dos recursos didático-pedagógicos presentes no cotidiano.
Deve -se buscar, então, a autonomia intelectual dos \textbf{alunos}, por meio de atividades planejadas pelo professor, para promover o uso de diversas habilidades de pensamento, como interpretar, analisar, sintetizar, classificar, relacionar e comparar, como possibilidade de desenvolver o processo de aprendizagem que \textbf{foca} uma formação \textbf{crítica} de estudantes e profissionais voltada para uma inserção consciente no mundo cultural e social.
Tal processo formativo deve favorecer a autonomia do educando, \textbf{despertando} -lhe a curiosidade, estimulando tomadas de decisões individuais e coletivas advindas das atividades essenciais da prática social e em contextos do estudante.
Pensar nessa transição requer propostas inovadoras que possam \textbf{auxiliar} o fazer pedagógico de cada unidade escolar a ter um olhar mais criterioso em relação às escolhas das temáticas das tendências em Educação Matemática e a sua aplicação como proposta \textbf{metodológica} para o ensino dessa \textbf{disciplina} ; deve -se identificar, no contexto de sua \textbf{aplicabilidade}, quais são as contribuições que ela proporciona no processo de promoção da cidadania e inclusão social, por meio do desenvolvimento de \textbf{competências} e habilidades que permitam aprender e continuar aprendendo, para compreender, questionar, interagir e tomar decisões.
O currículo por \textbf{competência} em Matemática que está sendo abordado neste material não é recente.
Desde a década de 1990, ele vem se fortalecendo no campo pedagógico brasileiro.
As \textbf{competências} são definidas neste documento como “ a \textbf{capacidade} do sujeito de \textbf{mobilizar} saberes, conhecimentos, habilidades e atitudes para resolver \textbf{problemas} e tomar decisões adequadas ” ( ZABALA, 1998 ) ; além disso, podem ser entendidas como \textbf{capacidade} de \textbf{mobilizar} recursos intelectuais/cognitivos para solucionar situações com pertinência ( GENTILE ; BENCINI, 2005 ).
Já para Perrenoud ( 1999 ), as \textbf{competências} elementares a serem desenvolvidas na escola têm relação com os programas escolares e com os saberes de natureza disciplinar, \textbf{exigindo} noções e conhecimentos de Matemática, Língua Portuguesa, Artes, Educação Física, Língua \textbf{Inglesa}, Ciências, História e Geografia.
Por outro \textbf{lado}, as habilidades estão relacionadas ao saber fazer, em uma dimensão mais técnica, e são necessárias para a consolidação de uma \textbf{competência}.
Nesse sentido, os eixos temáticos presentes na organização curricular por \textbf{competência} desta proposta aspiram a ser orientadores da formação de \textbf{competências} e habilidades, além de realizar aproximações com os objetos de conhecimentos referenciais para a formação dos estudantes em cada nível do Ensino Fundamental, colaborando com a organização \textbf{conceitual} e prática do que se considera essencial nas escolhas pedagógicas para cada ano.
O \textbf{fundamento} do “ Currículo por \textbf{Competências} ” é a redefinição do sentido dos conteúdos de ensino, de modo a atribuir sentido prático aos saberes escolares.
Vivemos hoje na era da \textbf{tecnologia} e da informação ; nunca se produziu e se consumiu \textbf{tanto} conteúdo na história da \textbf{humanidade}, em todos os níveis e áreas da sociedade.
Isso se deve à facilidade que temos em acessar essas informações e conteúdos, principalmente depois do \textbf{surgimento} e da expansão da internet.
Várias ferramentas tecnológicas estão disponíveis para as escolas adaptarem -se ao mundo \textbf{moderno} e incorporarem novos \textbf{métodos} de ensino que possam melhorar o processo de ensino e aprendizagem e a prática da \textbf{interdisciplinaridade}.
Uma das ferramentas que vem apresentando destaque é a robótica educacional, que \textbf{desperta} o interesse dos \textbf{alunos}, uma vez que eles podem explorar novas ideias e descobrir novos \textbf{caminhos} na aplicação de conceitos adquiridos em sala de \textbf{aula} e na resolução de \textbf{problemas}, desenvolvendo a \textbf{capacidade} de elaborar hipóteses, investigar soluções, estabelecer relações e tirar conclusões.
Nesse cenário, a escola teve que ( ou deve ) mudar seu \textbf{posicionamento}.
Antes dessa revolução da informação em nossa sociedade, a escola era tida como responsável pela disseminação de conteúdos.
Isso já não faz mais sentido, uma vez que os \textbf{alunos} têm acesso aos conteúdos, independentemente da escola, podendo, ainda, visualizá -los e consumi -los na quantidade, velocidade e no momento que desejarem.
Portanto, a escola deve \textbf{focar} seu trabalho em \textbf{competências} e habilidades para preparar o jovem para lidar com situações de seu cotidiano e ser capaz de resolver \textbf{problemas} reais.
Essa postura demonstra ainda alinhamento com as tendências educacionais que enfatizam a importância de colocar o \textbf{aluno} como protagonista, sendo um agente ativo em seu processo de ensino e aprendizagem, por meio, por exemplo, de atividades educativas extraclasses.
Nesse sentido, o conjunto de \textbf{competências} e habilidades básicas que transitam pelo direito de aprendizagem construído a partir da prática, argumentação, reflexão e produção de conhecimento.
O DCRB traz na sua proposta um olhar para a transição entre as \textbf{competências} gerais, \textbf{competências} específicas e habilidades que possa criar elos que ajudarão a desenvolver e substanciar, no âmbito pedagógico, os direitos de aprendizagem e desenvolvimento para resolver \textbf{demandas} \textbf{complexas} da vida cotidiana, do pleno exercício da cidadania e do mundo do trabalho.
As \textbf{competências} gerais apontam \textbf{caminhos} do aprendizado, avançam no campo das \textbf{socioemocionais} e norteiam o trabalho das escolas e dos professores em todos os anos e componentes curriculares, desenvolvendo, efetivamente, conhecimentos, habilidades e atitudes que vão se conectando com as habilidades referentes às áreas do conhecimento, \textbf{impactando} não apenas o currículo, mas processos de ensino e aprendizagem, gestão, formação de professores e avaliação.
Transição entre as \textbf{Competências} Gerais, \textbf{Competências} Específicas da Matemática e as Habilidades Fonte : Elaboração dos redatores do Componente Curricular de Matemática ( 2018 ). \textbf{COMPETÊNCIAS} GERAIS 1.
Valorizar e utilizar os conhecimentos historicamente construídos sobre o mundo físico, social, cultural e \textbf{digital} para entender e explicar a realidade, continuar aprendendo e colaborar para a construção de uma sociedade justa, democrática e inclusiva. \textbf{COMPETÊNCIAS} ESPECÍFICAS DE MATEMÁTICA 1.
Reconhecer que a Matemática é uma \textbf{ciência} humana, fruto das necessidades e preocupações de diferentes culturas, em diferentes momentos históricos, e é uma \textbf{ciência} viva, que contribui para solucionar \textbf{problemas} científicos e tecnológicos e para alicerçar descobertas e construções, inclusive com impactos no mundo do trabalho.
HABILIDADES 5.
Utilizar processos e ferramentas matemáticas, inclusive \textbf{tecnologias} \textbf{digitais} disponíveis, para modelar e resolver \textbf{problemas} cotidianos, sociais e de outras áreas de conhecimento, validando estratégias e resultados.
Dessa forma, a escola deve pensar práticas pedagógicas com o objetivo de desenvolver as habilidades dos estudantes.
A evolução das \textbf{competências} será fruto da \textbf{mobilização} dessas habilidades com o objetivo de resolver \textbf{problemas} e desafios.
Os conteúdos matemáticos recomendados por grande parte dos educadores matemáticos para a Educação Básica abrangem o desenvolvimento de habilidades de resolução de \textbf{problemas} envolvendo as operações fundamentais da Matemática, o \textbf{raciocínio} algébrico, o estabelecimento de relações, o reconhecimento de proporcionalidades e várias outras.
Essas habilidades são importantes, não somente para a trajetória escolar, mas para o próprio cotidiano da vida \textbf{moderna}.
No Ensino Fundamental, a escola precisa potencializar o estudante para entender como a Matemática é aplicada em diferentes situações, dentro e fora da escola.
Na \textbf{aula}, o contexto pode ser puramente matemático, ou seja, não é necessário que a questão apresentada seja referente a um fato cotidiano.
O importante é que os procedimentos sejam inseridos em uma rede de significados mais ampla na qual o \textbf{foco} não seja o cálculo em si, mas as relações que ele permite estabelecer entre os diversos conhecimentos que o \textbf{aluno} já tem.
Nossa expectativa é oferecer condições para que o \textbf{aluno} compreenda a Matemática em diferentes situações.
Evidencia -se, assim, a importância do conhecimento matemático como linguagem que, no diálogo com outros conhecimentos, amplia a \textbf{compreensão} do homem em relação ao mundo físico e social, aspecto que permite a resolução de situações-problema e a transformação da realidade.
Ao tratar da Matemática como componente curricular, este documento propõe cinco unidades temáticas correlacionadas, que orientam a formulação de habilidades a serem desenvolvidas ao longo do Ensino Fundamental : Números, Álgebra, Geometria, Grandezas e Medidas, Estatística e Probabilidade, as quais organizam os objetos de conhecimento ( conteúdos, conceitos e processos ) relacionados às suas respectivas habilidades ( aprendizagens essenciais que devem ser asseguradas aos estudantes nos diferentes contextos escolares ).
As expectativas de aprendizagens aumentam a cada nova etapa, bem como as habilidades que se \textbf{espera} desenvolver a partir do conhecimento construído em sala de \textbf{aula}.

% matched lemmas: abordagem, aluno, ancorar, análise, aplicabilidade, aprendiz, aprofundamento, aproximar, argumentar, artístico-cultural, ator, atual, aula, autoria, auxiliar, basear, caminho, capacidade, ciência, competência, complexo, compreensão, conceitual, consenso, contemporâneo, contextualização, crítico, demanda, depender, despertar, digital, disciplina, ensino-aprendizagem, esperar, exigir, figurar, final, focar, foco, fundamento, humanidade, impactar, implicar, imprescindível, inglês, inter-relacionar, interdisciplinaridade, introdutório, intuito, inúmero, lado, letramento, libra, metodológico, mobilizar, mobilização, moderno, método, posicionamento, problema, raciocínio, reflexivo, relevante, simples, socioambiental, socioemocional, surgimento, século, tanto, tecnologia, transversal, vista, visual-motor
\end{textsample}
